\documentclass[a4paper]{book}

% Load the VUB package.
% This has many options, please read the documentation at
% https://gitlab.com/rubdos/texlive-vub
\usepackage{vub}
\usepackage{emptypage}
\usepackage[dutch, british]{babel}
\usepackage{csquotes}
\usepackage{hyperref}
\usepackage{cleveref}

\setlength{\parskip}{\baselineskip}
\setlength{\parindent}{0pt} % Optional: removes paragraph first-line indent



% Some highly suggested packages, please read their manuals.
%\usepackage[natbib,style=apa]{biblatex}
%\addbibresource{bibliography.bib}

% LTeX: dictionary+=Haai
% LTeX: dictionary+=Gérard
% LTeX: dictionary+=Lichtert
% LTeX: dictionary+=Bjarno
% LTeX: dictionary+=Oeyen
% LTeX: dictionary+=Wolfgang
% LTeX: dictionary+=De Meuter
% LTeX: dictionary+=MCUs
% LTeX: dictionary+=Tierless

% LTeX: enabled=false
\title{Programming Robots with higher-Order Reactive programming}
\pretitle{\flushleft{Graduation thesis submitted in partial fulfilment of the
requirements for the degree of de
Ingenieurswetenschappen: Computerwetenschappen}}
\author{Gérard Lichtert}
\date{2025-2026}
\promotors{Promotor: Prof.\ Dr.\ Wolfgang De Meuter\\
Advisor: Bjarno Oeyen}
\faculty{Sciences and Bio-engineering Sciences} % Note: without the word "Faculty"!

\begin{document}
\frontmatter
\maketitle%

% Oftentimes, you need add a second language title.
\title{Robots programmeren met hogere-orde reactief programmeren}
\pretitle{\flushleft{Proefschrift ingediend met het oog op het behalen van de graad van Master of Science in de Ingenieurswetenschappen: Computerwetenschappen}}%
\date{2025-2026}%
\faculty{Wetenschappen en Bio-ingenieurswetenschappen}%
\maketitle%

\raggedright
% LTeX: enabled=true
\chapter*{Abstract}

\textbf{Reactive programming} (RP) is a declarative programming paradigm based on
data streams and the propagation of change. In a nutshell, if we have two
variables x and y and another variable z that sums both x and y, it should be sufficient to
update either x or y and z should be updated automatically.
\par
There are several approaches to achieve reactive programming. One approach is to
create a dedicated language. In this case the reactivity will be native to the
language. Another option is to extend a language with reactive features, like
%example

\tableofcontents

\mainmatter

\chapter{Introduction}
Modern software is written in such a way that it is event-driven, or in other
words, reactive. This means that applications reacts to an incoming event
by creating subsequent events that arise as a reaction to the incoming event.
Such an application is often called a reactive system. Reactive Systems are
often implemented with callbacks or by implementing the Observer Pattern.
However, this does not come without its drawbacks as this may lead to
\enquote{inversion of control} or \enquote{callback hell}. To alleviate this we
introduce \textbf{Reactive Programming}. A programming paradigm that allows the
user to write reactive systems in a declarative way.
There are different ways to achieve Reactive Programming. One way is to create a
reactive programming language, making the language itself reactive. Another is
to extend a language with reactive features, either by introducing native
reactivity or through a library. We will be looking at the former way and more
specifically into the higher order reactive programming language \textbf{Haai}.
\par
Additionally, we want to delve into the world of electronics.
From maintaining your fridge's temperature to controlling a drone's
altitude, electronics are deeply embedded in everyone's day-to-day life. Writing
software for these electronics is also a challenge, mainly due to the fact that
using peripherals require using their associated libraries as well as using a
low-level programming language such as C/C++ or Rust.
\section{Contributions}
\section{Overview of this Dissertation}
\section{Approach}
\chapter{State of the Art}
\section{Reactive Programming}

\section{RP languages targeting MCUs}
\section{Systems to program robots}
\section{Tierless Programming}
\section{Problem Statement}
\chapter{$\mu$-Remus}
\section{Haai \& Remus}
\section{ROS \& $\mu$-ROS}
\section{Mapping RP to ROS \& $\mu$-ROS}
\section{$\mu$-Remus: A C implementation}
\chapter{Programming Robots with Reactors}
\section{Linking ROS and $\mu$-ROS}
\chapter{Tierless Haai}
\section{Haai to CHaai}
\chapter{Evaluation}
\section{Evaluating Haai distribution}
\section{Evaluating Haai to CHaai}
\section{Evaluating CHaai on $\mu$-Remus}
\chapter{Conclusion \& Future Work}
\backmatter%







%\printbibliography%
\end{document}
